\documentclass[10pt,a4paper]{amsart}
\usepackage{amsmath,amssymb,amsthm,graphicx}
\usepackage{verbatim}
\newenvironment{reference}{\comment}{\endcomment}
%\newenvironment{reference}{}{}
\newenvironment{slogan}{\comment}{\endcomment}
\newenvironment{history}{\comment}{\endcomment}
% For commutative diagrams we use Xy-pic
\usepackage[all]{xy}

% We use 2cell for 2-commutative diagrams.
\xyoption{2cell}
% Theorem environments.
%
\theoremstyle{plain}
\newtheorem{theorem}[subsection]{Theorem}
\newtheorem{proposition}[subsection]{Proposition}
\newtheorem{lemma}[subsection]{Lemma}

\theoremstyle{definition}
\newtheorem{definition}[subsection]{Definition}
\newtheorem{example}[subsection]{Example}
\newtheorem{exercise}[subsection]{Exercise}
\newtheorem{situation}[subsection]{Situation}

\theoremstyle{remark}
\newtheorem{remark}[subsection]{Remark}
\newtheorem{remarks}[subsection]{Remarks}

\numberwithin{equation}{subsection}

% Macros
%
\def\lim{\mathop{\mathrm{lim}}\nolimits}
\def\colim{\mathop{\mathrm{colim}}\nolimits}
\def\Spec{\mathop{\mathrm{Spec}}}
\def\Hom{\mathop{\mathrm{Hom}}\nolimits}
\def\Ext{\mathop{\mathrm{Ext}}\nolimits}
\def\SheafHom{\mathop{\mathcal{H}\!\mathit{om}}\nolimits}
\def\SheafExt{\mathop{\mathcal{E}\!\mathit{xt}}\nolimits}
\def\Sch{\mathit{Sch}}
\def\Mor{\mathop{\mathrm{Mor}}\nolimits}
\def\Ob{\mathop{\mathrm{Ob}}\nolimits}
\def\Sh{\mathop{\mathit{Sh}}\nolimits}
\def\NL{\mathop{N\!L}\nolimits}
\def\CH{\mathop{\mathrm{CH}}\nolimits}
\def\proetale{{pro\text{-}\acute{e}tale}}
\def\etale{{\acute{e}tale}}
\def\QCoh{\mathit{QCoh}}
\def\Ker{\mathop{\mathrm{Ker}}}
\def\Im{\mathop{\mathrm{Im}}}
\def\Coker{\mathop{\mathrm{Coker}}}
\def\Coim{\mathop{\mathrm{Coim}}}

% Boxtimes
%
\DeclareMathSymbol{\boxtimes}{\mathbin}{AMSa}{"02}

%
% Macros for moduli stacks/spaces
%
\def\QCohstack{\mathcal{QC}\!\mathit{oh}}
\def\Cohstack{\mathcal{C}\!\mathit{oh}}
\def\Spacesstack{\mathcal{S}\!\mathit{paces}}
\def\Quotfunctor{\mathrm{Quot}}
\def\Hilbfunctor{\mathrm{Hilb}}
\def\Curvesstack{\mathcal{C}\!\mathit{urves}}
\def\Polarizedstack{\mathcal{P}\!\mathit{olarized}}
\def\Complexesstack{\mathcal{C}\!\mathit{omplexes}}
% \Pic is the operator that assigns to X its picard group, usage \Pic(X)
% \Picardstack_{X/B} denotes the Picard stack of X over B
% \Picardfunctor_{X/B} denotes the Picard functor of X over B
\def\Pic{\mathop{\mathrm{Pic}}\nolimits}
\def\Picardstack{\mathcal{P}\!\mathit{ic}}
\def\Picardfunctor{\mathrm{Pic}}
\def\Deformationcategory{\mathcal{D}\!\mathit{ef}}
\title{The cokernel of a map from a module of finite type into a finitely presented module is finitely presented}
\author{The Stacks Project Authors}
\date{stacks-project:modules.tex:local-snapshot-20261012; GNU FDL 1.2 or later}
\begin{document}
\maketitle
\noindent\textit{Editorial note.} All delivered mathematical text is the author's own, copied from the pinned Stacks Project chapter source sources/stacks-modules/pinned-original-source.tex; no statement or proof step was rewritten, reordered or removed.\par\medskip
\noindent\textit{Editorial note.} Preamble provenance: the Stacks Project keeps one single shared preamble for every chapter; this card takes it directly from the staged sources/stacks-modules/preamble.tex as three verbatim spans (lines 8-12, 14-18 and 154-217), so no derived file was written and nothing was added or redefined. The wrapper class and the base packages are supplied by the assembler.\par\medskip
\noindent\textit{Editorial note.} Bibliography provenance: neither the counted statement nor its proof contains a citation (this lane rescans the whole span before writing the card), so the card carries no bibliography and no entry was reconstructed.\par\medskip
\noindent\textit{Editorial note.} Selection provenance: this card was opened under the file-based control contract (CONTROL.md), with the mechanical pre-card gate run on its statement and proof interval before the card was written.\par\medskip
\noindent\textit{Allowed standard prerequisites (editorial).}\begin{itemize}
\item The local description of a finitely presented module as a quotient of a finite direct sum of the structure sheaf by a submodule of finite type. Used at the start of the counted proof (source lines 1658-1661). Provides the free module and the finite type submodule in which the computation takes place.
\item The finiteness of the extension in a short exact sequence whose outer terms are of finite type, and the finiteness of images of finite type modules. Author lemma lemma-extension-finite-type, whose statement is reproduced verbatim as a context span of this card (source lines 1071-1081). Shows that the submodule generated by the image of the given map is again of finite type.
\end{itemize}\medskip
\begin{lemma}
\label{lemma-cokernel-finite-finite-presentation}
Let $(X,\mathcal{O}_X)$ be a ringed space. Let $\mathcal{F}$ be a
finitely presented $\mathcal{O}_X$-module. Let
$\varphi : \mathcal{G} \to \mathcal{F}$ be a morphism of
$\mathcal{O}_X$-modules. If $\mathcal{G}$ is finite type, then
$\Coker(\varphi)$ is finitely presented.
\end{lemma}

\begin{lemma}
\label{lemma-extension-finite-type}
Let $X$ be a ringed space.
The image of a morphism of $\mathcal{O}_X$-modules of finite
type is of finite type.
Let
$0 \to \mathcal{F}_1 \to \mathcal{F}_2 \to \mathcal{F}_3 \to 0$
be a short exact sequence of $\mathcal{O}_X$-modules.
If $\mathcal{F}_1$ and $\mathcal{F}_3$ are of finite type,
so is $\mathcal{F}_2$.
\end{lemma}

\begin{proof}
Locally on $X$, we can write
$\mathcal{F} = \mathcal{O}_X^{\oplus n} / \mathcal{M}$, where
$\mathcal{M} \subset \mathcal{O}_X^{\oplus n}$ is a finite type
$\mathcal{O}_X$-submodule. Therefore
$\Im(\varphi) = \mathcal{N} / \mathcal{M}$, where
$\mathcal{N} \subset \mathcal{O}_X^{\oplus n}$ is a
$\mathcal{O}_X$-submodule containing $\mathcal{M}$.
The $\mathcal{O}_X$-module $\Im(\varphi)$ is of finite type
because $\mathcal{G} \to \Im(\varphi)$ is surjective and
$\mathcal{G}$ is finite type. By Lemma \ref{lemma-extension-finite-type},
$\mathcal{N}$ is finite type. Thus
$\Coker(\varphi) = \mathcal{O}_X^{\oplus n} / \mathcal{N}$
is finitely presented.
\end{proof}

\end{document}
